\chapter{Scope and Qualifications}

\section{Reader and outcome}

This is a print-first, self-paced language course for a learner fluent in
English, familiar with formal English grammar and the structure of Mass, but
beginning with no Latin grammar. Its controlling reader is the 1962 typical
edition of the \emph{Missale Romanum}; the Clementine Vulgate and other
ecclesiastical prose are supporting corpora. The intended outcome is functional
independence in grammatical analysis and direct reading, followed by controlled
and guided composition. Pronunciation, chant, ceremonial instruction, current
authorization to use the books, and a history of Latin are outside scope.

The course does not certify mastery, confer ecclesiastical approval, or qualify
a reader to settle a liturgical, theological, textual, or canonical question.
Knowledge of Latin improves access to sources but does not replace competence
in those disciplines.

\section{Text, orthography, and translation}

Received Missal wording is identified by liturgical location and is to be
checked against the 1962 typical-edition page images named in the research
records. Biblical comparison uses the identified 1914 Hetzenauer Clementine
Vulgate unless a local label names another witness. Searchable presentations,
OCR, parsers, and dictionary interfaces are finding aids rather than controlling
texts.

Running Latin normally follows the course house typography: consonantal
\Latin{i}, the ligature \Latin{æ}, and no macrons. A source passage may therefore
be normalized typographically without claiming a diplomatic transcription.
Textual-comparison exercises preserve or explicitly label witness spelling and
normalization. Uncited short drills, transformations, and model compositions
are project-composed. Source-dependent liturgical, scriptural, or historical
passages are received wording or are explicitly recorded in the passage
inventory as excerpts, adaptations, abridgments, or composites. English
translations and glosses are project translations unless an identified human
translation is named locally.

\section{Completeness, sources, and review}

“Complete curriculum” means that the course owner preserves the supplied
roadmap, reference grammar, four teaching stages, worksheet banks, practice
forms, parallel mastery assessments, and their keys. Those sources are
published through prerequisite-ordered module and companion packets rather
than one omnibus file. A normal learning packet prints two selected worksheet
variants and identifies the remaining two as optional remediation or delayed
work. Completeness does not mean a complete
historical grammar, lexicon, patristic reader, critical edition, or inventory of
all Ecclesiastical Latin. Rare vocabulary and advanced textual questions still
require appropriate editions, lexica, and specialist judgment.

The working manuscript was supplied for integration from a separate course
project. Its generation provenance was not recorded in that project. Triptych's
integration records preserve the supplied source identity, content map,
corrections, and known limits without attributing the donor text to an
unverified author or model. Internal grammar, source, key, and production checks
are recorded in the adjacent research files. No independent Latinist,
liturgist, pedagogue, theologian, or ecclesiastical reviewer is claimed.

\section{Rights and use}

The integration request identified and authorized work on the donor corpus,
but it did not independently establish authorship, ownership, or a distribution
license for that corpus; no such record was present in the donor project.
Scripture, Missal wording and rubrics, historical ecclesiastical text, and
third-party reference works also remain outside the Triptych license. Their
inclusion does not promise that they may be extracted or redistributed
separately. The source records identify known public-domain bases and unresolved
work-specific distribution questions. This snapshot remains outside public
release unless and until its exact bytes and rights record receive separate
authorization.
